\section{Method and Analysis \tagO}\label{sec:method}

\begin{figure}[t]
\centering
% Pipeline figure, shared by the report (sections/04_method.tex) and the poster.
% Needs: TikZ libraries positioning, arrows.meta, calc; colour cOurs.
% Sizes are absolute (mm), so the font is fixed at 7 pt; the poster scales the whole picture.
% Grid-aligned: both encoder lanes have identical widths, all arrows are horizontal or vertical.
\begin{tikzpicture}[
  font=\fontsize{7}{8.5}\selectfont,
  enc/.style={draw=black!45, fill=black!5, text width=29mm, align=center,
              inner sep=1.2mm, minimum height=8mm},
  lora/.style={draw=cOurs, fill=cOurs!12, text width=29mm, align=center,
               inner sep=1.2mm, minimum height=4.5mm},
  op/.style={draw=black!60, fill=white, align=center, inner sep=1.2mm, minimum height=9mm},
  loss/.style={draw=cOurs, fill=cOurs!12, align=center, inner sep=1.2mm, minimum height=9mm},
  io/.style={align=center, inner sep=0.5mm},
  arr/.style={-{Stealth[length=1.6mm]}, draw=black!70}]
% ---- image lane (top) and text lane (bottom), same x positions
\node[enc]  (vit) at (0, 7.5mm)  {ViT-B/16 image encoder\\12 layers, frozen};
\node[lora, below=0pt of vit]     (lv)  {+ LoRA on $W_{q,k,v,o}$};
\node[enc]  (tr)  at (0,-11mm)   {Text transformer\\12 layers, frozen};
\node[lora, below=0pt of tr]      (lt)  {+ LoRA on $W_{q,k,v,o}$};
\node[io, left=3mm of vit] (img) {person image\\$384{\times}128$};
\node[io, left=3mm of tr]  (txt) {description\\$\le$77 tokens};
\node[io, right=3mm of vit] (v) {$\mathbf{v}_i$};
\node[io, right=3mm of tr]  (t) {$\mathbf{t}_j$};
% ---- similarity and loss, centred between the two lanes
\path (v) -- (t) coordinate[pos=0.5] (mid);
\node[op]   (S) at ($(mid)+(10mm,0)$) {$B{\times}B$\\cosine$/\tau$};
\node[loss, right=4mm of S] (L) {$\mathcal{L}_{\mathrm{sdm}}$\\only loss};
\node[io, above=4mm of L] (y) {identities $y$};
% ---- arrows (orthogonal only)
\draw[arr] (img) -- (vit);
\draw[arr] (txt) -- (tr);
\draw[arr] (vit) -- (v);
\draw[arr] (tr)  -- (t);
\draw[arr] (v) -| (S.north);
\draw[arr] (t) -| (S.south);
\draw[arr] (S) -- (L);
\draw[arr] (y) -- (L);
\end{tikzpicture}

\caption{The pipeline. Grey: frozen CLIP ViT-B/16 (\nBackboneM\,M parameters). Blue: the only trainable parameters, \nLoraModules{} LoRA modules of rank 4 (\nTrainable{} parameters), trained with SDM alone.}
\label{fig:pipeline}
\end{figure}

\subsection{LoRA-tuned CLIP}\label{sec:lora}

Every linear projection inside multi-head attention, the query, key, value and output maps $W\in\mathbb{R}^{d\times d}$, is replaced by
\begin{equation}
W'\mathbf{x}=W\mathbf{x}+\frac{\alpha}{r}\,BA\,\mathbf{x},\qquad A\in\mathbb{R}^{r\times d},\;B\in\mathbb{R}^{d\times r},\label{eq:lora}
\end{equation}
with $W$ frozen, $r=4$ and $\alpha=8$. $A$ is initialised Kaiming-uniform and $B$ to zero, so training starts exactly from pre-trained CLIP. Applying this to both encoders gives $12\times4\times2=\nLoraModules$ modules (\cref{fig:pipeline}). Each module holds $r(d+d)$ parameters, so the image encoder ($d=768$) contributes $48\cdot4\cdot1536=294{,}912$ and the text encoder ($d=512$) $48\cdot4\cdot1024=196{,}608$, a total of \nTrainable, or \trainPct\% of the \nTotalM\,M parameters of the adapted model. Everything else stays frozen, including the embeddings, layer norms, feed-forward blocks, projection heads and the temperature. Since $BA$ can be added to $W$ once training ends, inference costs exactly as much as CLIP's.

\textbf{Objective.} We train with $\mathcal{L}_{\mathrm{sdm}}$ of \cref{eq:sdm} and nothing else. Besides isolating the loss, this keeps the trainable part small: IRRA's ID classifier alone would add $512\times3701$ weights plus biases, \idHeadM\,M parameters or nearly four times our LoRA budget, and the IRR encoder another 13.66\,M. SDM has no parameters and needs the labels only through the relation $y_i=y_j$.

\subsection{What SDM computes with identity labels}\label{sec:analysis}

\textbf{How often does a batch contain a second positive?} SDM differs from InfoNCE only in rows of $q$ that have more than one non-zero entry. RSTPReid's training split contains \nExamples{} image--caption examples, ten per identity (five images, two captions each). With uniform random batches of $B=\batchSize$, as in IRRA's default and our implementation, the probability that a given row has at least one other example of its identity in the batch is
\begin{equation}
1-\prod_{k=0}^{B-2}\frac{N-10-k}{N-1-k}=\posFrac\%,\qquad N=37{,}010.\label{eq:posfrac}
\end{equation}
In \oneHotFrac\% of rows, therefore, $q_{i,\cdot}$ is the one-hot vector $\mathbf{e}_i$, exactly InfoNCE's target. Any difference between SDM and InfoNCE in this regime must come from the form of the loss, not from the targets.

\textbf{SDM with a one-hot target is not InfoNCE.} Substituting $q_{i,\cdot}=\mathbf{e}_i$ into one row of \cref{eq:sdm} and writing $\lambda=\log(1/\epsilon)$ gives
\begin{align}
\ell^{\mathrm{sdm}}_i&=\sum_{j}p_{i,j}\log p_{i,j}+\lambda\sum_{j\neq i}p_{i,j}-p_{i,i}\log(1+\epsilon)\nonumber\\
&\approx \underbrace{\lambda\,(1-p_{i,i})}_{\text{scaled MAE}}\;-\;\underbrace{H(p_{i,\cdot})}_{\text{entropy bonus}},\label{eq:decomp}
\end{align}
with $\lambda=\logInvEps$, whereas InfoNCE's row loss is $\ell^{\mathrm{nce}}_i=-\log p_{i,i}$. The first term is the mean absolute error between $p_{i,\cdot}$ and $\mathbf{e}_i$, $\tfrac12\lVert p_{i,\cdot}-\mathbf{e}_i\rVert_1=1-p_{i,i}$, scaled by $\lambda$; MAE is known to be robust to label noise because it is a symmetric, bounded loss~\cite{ghosh2017robustlossfunctionslabel}. The second term rewards high-entropy predictions, which is the confidence penalty of Pereyra~\textit{et al.}~\cite{pereyra2017regularizingneuralnetworkspenalizing} with weight $1/\lambda\approx0.05$ relative to the first. The two losses share a minimiser, $p_{i,\cdot}=\mathbf{e}_i$, but weight the examples very differently.

\textbf{Gradients.} For logits $z_{i,j}=\mathbf{v}_i^{\top}\mathbf{t}_j/\tau$ and $\tilde q=q+\epsilon$, the exact gradient of one SDM row is
\begin{equation}
\frac{\partial\ell^{\mathrm{sdm}}_i}{\partial z_{i,j}}=p_{i,j}\big(\log p_{i,j}-\log\tilde q_{i,j}-\ell^{\mathrm{sdm}}_i\big),\label{eq:grad}
\end{equation}
while InfoNCE gives $p_{i,j}-\mathbf{1}[j=i]$. On the matching pair, InfoNCE pushes with strength $1-p_{i,i}$, which is largest for the pairs the model ranks worst, while the dominant MAE term of \cref{eq:decomp} pushes with $\lambda\,p_{i,i}(1-p_{i,i})$, which vanishes as $p_{i,i}\to0$. SDM thus concentrates on pairs the model already half-recognises and largely ignores pairs it cannot match, the opposite of InfoNCE's emphasis on the hardest pairs. In TBPS such pairs are often ambiguous rather than informative: two people in near-identical clothing, or a caption that omits the only distinguishing detail.

\textbf{A control loss that isolates the KL direction.} To test whether this, rather than the identity-aware target, explains SDM's advantage, we introduce multi-positive cross-entropy,
\begin{align}
\mathcal{L}^{\mathrm{i2t}}_{\mathrm{mpce}}&=-\frac1B\sum_{i}\sum_{j}q_{i,j}\log p_{i,j}\nonumber\\
&=\frac1B\sum_i\Big[\KL(q_{i,\cdot}\Vert p_{i,\cdot})+H(q_{i,\cdot})\Big],\label{eq:mpce}
\end{align}
summed over both directions as for SDM. MPCE uses exactly SDM's targets but the \emph{forward} KL, and $H(q)$ does not depend on the model. The three losses then form two controlled contrasts: MPCE versus InfoNCE changes only the target (identity-aware or diagonal), and SDM versus MPCE changes only the KL direction. In the one-hot regime MPCE coincides with InfoNCE up to the factor $2$ of \cref{eq:nce}, and Adam-type optimisers are invariant to a constant rescaling of the loss (up to their $\epsilon$ and gradient clipping). Our hypothesis therefore predicts that MPCE will \emph{behave like InfoNCE}, not like SDM. All three identities (\cref{eq:decomp}, MPCE $=2\times$InfoNCE for distinct identities, and SDM $\neq$ InfoNCE) are checked numerically by unit tests in our code.
